\begin{register}{H}{SDHOST\_CTRL\_REG}{0x{}0000}\label{regdesc:CTRLREG}
\regfield{(reserved)}{7}{25}{{0x{}00}}%
\regfield{(reserved)}{1}{24}{1}%
\regfield{(reserved)}{12}{12}{{0x{}000}}%
\regfield{SDHOST\_CEATA\_DEVICE\_INTERRUPT\_STATUS}{1}{11}{0}%
\regfield{SDHOST\_SEND\_AUTO\_STOP\_CCSD}{1}{10}{0}%
\regfield{SDHOST\_SEND\_CCSD}{1}{9}{0}%
\regfield{SDHOST\_ABORT\_READ\_DATA}{1}{8}{0}%
\regfield{SDHOST\_SEND\_IRQ\_RESPONSE}{1}{7}{0}%
\regfield{SDHOST\_READ\_WAIT}{1}{6}{0}%
\regfield{(reserved)}{1}{5}{0}%
\regfield{SDHOST\_INT\_ENABLE}{1}{4}{0}%
\regfield{(reserved)}{1}{3}{0}%
\regfield{SDHOST\_DMA\_RESET}{1}{2}{0}%
\regfield{SDHOST\_FIFO\_RESET}{1}{1}{0}%
\regfield{SDHOST\_CONTROLLER\_RESET}{1}{0}{0}%
\reglabel{Reset}\regnewline%
\begin{regdesc}\begin{reglist}
\item [SDHOST\_CEATA\_DEVICE\_INTERRUPT\_STATUS] 软件应在 CE-ATA 设备上电复位或其他复位后写此位。CE-ATA 设备的中断通常被禁能 (nIEN = 1)。如果主机使能中断，则软件应将此位置为 1。(R/W)
\item [SDHOST\_SEND\_AUTO\_STOP\_CCSD] 始终将 SDHOST\_SEND\_AUTO\_STOP\_CCSD 和 SDHOST\_SEND\_CCSD 一起置位，不能分开单独置位。置位时， SD/MMC 自动发送内部生成的 STOP 命令 (CMD12) 给 CE-ATA 设备。发送 STOP 命令后，SDHOST\_RINTSTS\_REG 里的 Auto Command Done (ACD) 位被置为 1，如果 ACD 中断没有屏蔽，则用于主机的中断将会生成。发送 Command Completion Signal Disable (CCSD) 后，SD/MMC 自动清零 SDHOST\_SEND\_AUTO\_STOP\_CCSD 位。(R/W)
\item [SDHOST\_SEND\_CCSD] 置位时，SD/MMC 发送 CCSD 给 CE-ATA 设备。只有在当前命令等待 CCS 并且 CE-ATA 设备的中断使能时软件才会将此位置位。一旦 CCSD 模式发送给设备，SD/MMC 就会自动清零 SDHOST\_SEND\_CCSD 位。如果 ACD 中断没有屏蔽，则寄存器 SDHOST\_RINTSTS\_REG 里的 Command Done (CD) 位会被置为 1，且主机的中断将会生成。\\
说明：一旦 SDHOST\_SEND\_CCSD 被置位，需要两个卡时钟周期来驱动 CMD 线上的 CCSD。因此，即使设备已经发送 CCS 信号，在边界条件内 CCSD 信号可能会发送给 CE-ATA 设备。(R/W)
\item [SDHOST\_ABORT\_READ\_DATA] 在读操作期间，挂起命令发送后，软件会轮询卡并找出挂起事件是从什么时候开始的。一旦挂起事件开始，软件会复位数据状态机，此时状态机正在等待下一个数据块。当数据状态机复位为空闲状态时，此位自动清零。(R/W)
\item [SDHOST\_SEND\_IRQ\_RESPONSE] 响应发送后此位自动清零。为了等待 MMC 卡发送的中断，主机发送 CMD40，然后等待 MMC 卡的中断响应。同时，如果主机想要 SD/MMC 退出等待中断的状态，则会将此位置 1，此时 SD/MMC 命令状态机发送 CMD40 响应并返回空闲状态。(R/W)
\item [见下页...]
\end{reglist}\end{regdesc}
\end{register}

\addtocounter{Regfloat}{-1}
\begin{register}{H}{SDHOST\_CTRL\_REG}{0x{}0000}
\begin{regdesc}\begin{reglist}
\item [接上页...]
\item [SDHOST\_READ\_WAIT] 发送读操作等待给 SDIO 卡。(R/W)
\item [SDHOST\_INT\_ENABLE] 全局中断使能/禁能位。0：禁能；1：使能。(R/W)
\item [SDHOST\_DMA\_RESET] 要复位 DMA 接口，软件应将此位置为 1。两个 AHB 时钟后此位自动清零。(R/W)
\item [SDHOST\_FIFO\_RESET] 要复位 FIFO，软件应将此位置为 1。复位操作结束后自动清零。\\说明：清零后，再过两个系统时钟周期和同步延迟（两个卡时钟周期），FIFO 指针将会退出复位。(R/W)
\item [SDHOST\_CONTROLLER\_RESET] 要复位控制器，软件应将此位置为 1。两个 AHB 时钟周期和两个 sdhost\_cclk\_in 时钟周期后此位自动清零。(R/W)
\end{reglist}\end{regdesc}
\end{register}

\begin{register}{H}{SDHOST\_CLKDIV\_REG}{0x{}0008}\label{regdesc:CLKDIVREG}
\regfield{SDHOST\_CLK\_DIVIDER3}{8}{24}{{0x{}00}}%
\regfield{SDHOST\_CLK\_DIVIDER2}{8}{16}{{0x{}00}}%
\regfield{SDHOST\_CLK\_DIVIDER1}{8}{8}{{0x{}00}}%
\regfield{SDHOST\_CLK\_DIVIDER0}{8}{0}{{0x{}00}}%
\reglabel{Reset}\regnewline%
\begin{regdesc}\begin{reglist}
\item [SDHOST\_CLK\_DIVIDER\textcolor{red}{\it{m}}]Clock divider-\textcolor{red}{\it{m}} 的值。时钟分频系数为 2*\textcolor{red}{\it{n}}，\textcolor{red}{\it{n}}=0 旁路分频器（分频系数为 1）。例如，值为 1 代表分频系数为 2*1 = 2，值为 0x{}FF 代表分频系数为 2*255 = 510，以此类推。m 的取值范围为 0 至 3。(R/W)
\end{reglist}\end{regdesc}
\end{register}

\begin{register}{H}{SDHOST\_CLKSRC\_REG}{0x{}000C}\label{regdesc:CLKSRCREG}
\regfield{(reserved)}{28}{4}{{0x{}0000000}}%
\regfield{SDHOST\_CLKSRC\_REG}{4}{0}{{0x{}0}}%
\reglabel{Reset}\regnewline%
\begin{regdesc}\begin{reglist}
\item [SDHOST\_CLKSRC\_REG] 时钟分频源可以支持 2 个 SD 卡。每个卡分配有两个位。例如，bit[1:0] 分配给卡 0，bit[3:2] 分配给卡 1。 卡 0 根据位值将时钟分频器 [0:3] 的输出信号传输给 cclk\_out[1:0] 管脚。(R/W)\\
00：时钟分频器 0；\\
01：时钟分频器 1；\\
10：时钟分频器 2；\\
11：时钟分频器 3。
\end{reglist}\end{regdesc}
\end{register}


\begin{register}{H}{SDHOST\_CLKENA\_REG}{0x{}0010}\label{regdesc:CLKENAREG}

\regfield{(reserved)}{14}{18}{{0x{}0000}}%
\regfield{SDHOST\_LP\_ENABEL}{2}{16}{{0x{}0}}%
\regfield{(reserved)}{14}{2}{{0x{}0000}}%
\regfield{SDHOST\_CCLK\_ENABEL}{2}{0}{{0x{}0}}%
\reglabel{Reset}\regnewline%
\begin{regdesc}\begin{reglist}
\item [SDHOST\_LP\_ENABLE] 当卡处于 IDLE 状态时，把时钟停掉。每个卡分配有一个位。(R/W)\\
0：时钟禁用；\\
1：时钟使能。
\item [SDHOST\_CCLK\_ENABLE] 时钟使能控制可支持 2 个 SD 卡时钟和一个 MMC 卡时钟。每个卡分配有一个位。(R/W)\\
0：时钟禁用；\\
1：时钟使能。
\end{reglist}\end{regdesc}
\end{register}


\begin{register}{H}{SDHOST\_TMOUT\_REG}{0x{}0014}\label{regdesc:TMOUTREG}
\regfield{SDHOST\_DATA\_TIMEOUT}{24}{8}{{0x{}FFFFFF}}%
\regfield{SDHOST\_RESPONSE\_TIMEOUT}{8}{0}{{0x{}40}}%
\reglabel{Reset}\regnewline%
\begin{regdesc}\begin{reglist}
\item [SDHOST\_DATA\_TIMEOUT] 此位用于设置卡数据读取超时的值，还可以用来设置主机超时的定时器的值。只有当卡时钟停止后超时计数器才开始启动。此位的值以卡输出时钟数来表示，即被选取卡的 sdhost\_cclk\_out 时钟。(R/W)\\
说明：如果超时值是 100 ms 左右，则应该使用软件定时器。这种情况下，读数据超时中断应该被禁能。
\item [SDHOST\_RESPONSE\_TIMEOUT] 此位用于设置响应超时的值，以卡输出时钟数来表示，即 sdhost\_cclk\_out 时钟。(R/W)
\end{reglist}\end{regdesc}
\end{register}


\begin{register}{H}{SDHOST\_CTYPE\_REG}{0x{}0018}\label{regdesc:CTYPEREG}
\regfield{(reserved)}{14}{18}{{0x{}0000}}%
\regfield{SDHOST\_CARD\_WIDTH8}{2}{16}{{0x{}0}}%
\regfield{(reserved)}{14}{2}{{0x{}0000}}%
\regfield{SDHOST\_CARD\_WIDTH4}{2}{0}{{0x{}0}}%
\reglabel{Reset}\regnewline%
\begin{regdesc}\begin{reglist}
\item [SDHOST\_CARD\_WIDTH8] 每个卡一个位，表明卡是否处于 8-bit 模式。(R/W)\\
0：非 8-bit 模式；\\
1：8-bit 模式。\\
Bit[17:16] 分别对应卡 [1:0]。
\item [SDHOST\_CARD\_WIDTH4] 每个卡一个位，表明卡处于 1-bit 模式还是 4-bit 模式。(R/W)\\
0：1-bit 模式；\\
1：4-bit 模式。\\
Bit[1:0] 分别对应卡 [1:0]。
\end{reglist}\end{regdesc}
\end{register}


\begin{register}{H}{SDHOST\_BLKSIZ\_REG}{0x{}001C}\label{regdesc:BLKSIZREG}
\regfield{(reserved)}{16}{16}{{0x{}0000}}%
\regfield{SDHOST\_BLOCK\_SIZE}{16}{0}{{0x{}200}}%
\reglabel{Reset}\regnewline%
\begin{regdesc}\begin{reglist}
\item [SDHOST\_BLOCK\_SIZE] 块大小。(R/W)
\end{reglist}\end{regdesc}
\end{register}


\begin{register}{H}{SDHOST\_BYTCNT\_REG}{0x{}0020}\label{regdesc:BYTCNTREG}
\regfield{}{32}{0}{{0x{}200}}%
\reglabel{Reset}\regnewline%
\begin{regdesc}\begin{reglist}
\item [SDHOST\_BYTCNT\_REG] 表明传输的字节数。对于块的传输，值应为块大小的整数倍。对于未定义字节长度的数据传输，字节计数应该设置为 0。当字节计数为 0 时，主机应当明确发送停止/终止命令来停止数据传输。(R/W)
\end{reglist}\end{regdesc}
\end{register}


\begin{register}{H}{SDHOST\_INTMASK\_REG}{0x{}0024}\label{regdesc:INTMASKREG}
\regfield{(reserved)}{14}{18}{{0x{}0000}}%
\regfield{SDHOST\_SDIO\_INT\_MASK}{2}{16}{{0x{}0}}%
\regfield{SDHOST\_INT\_MASK}{16}{0}{{0x{}0000}}%
\reglabel{Reset}\regnewline%
\begin{regdesc}\begin{reglist}
\item [SDHOST\_SDIO\_INT\_MASK]SDIO 中断的屏蔽位，每个卡一个位。Bit[17:16] 分别对应卡 [1:0]。当被屏蔽时，SDIO 中断的检测被禁能。0 屏蔽中断，1 使能中断。(R/W)
\item [SDHOST\_INT\_MASK] 这些位用于屏蔽不想要的中断。0 屏蔽中断，1 使能中断。(R/W)\\
Bit 15 (EBE): 结束位错误/无 CRC 错误；\\
Bit 14 (ACD): 自动命令结束；\\
Bit 13 (SBE/BCI): 接收启动位错误；\\
Bit 12 (HLE): 硬件锁定写入错误\\
Bit 11 (FRUN): FIFO 空/满错误；\\
Bit 10 (HTO): 主机填充数据超时；\\
Bit 9 (DRTO): 数据读取超时；\\
Bit 8 (RTO): 响应超时； \\
Bit 7 (DCRC): 数据 CRC 错误； \\
Bit 6 (RCRC): 响应 CRC 错误； \\
Bit 5 (RXDR): 接收 FIFO 数据请求； \\
Bit 4 (TXDR): 发送 FIFO 数据请求； \\
Bit 3 (DTO): 数据传输结束； \\
Bit 2 (CD): 命令执行完毕； \\
Bit 1 (RE): 响应错误；\\
Bit 0 (CD): 卡检测。
\end{reglist}\end{regdesc}
\end{register}


\begin{register}{H}{SDHOST\_CMDARG\_REG}{0x{}0028}\label{regdesc:CMDARGREG}
\regfield{}{32}{0}{{0x{}00000000}}%
\reglabel{Reset}\regnewline%
\begin{regdesc}\begin{reglist}
\item [SDHOST\_CMDARG\_REG] 传递给卡的命令参数。(R/W)
\end{reglist}\end{regdesc}
\end{register}


\begin{register}{H}{SDHOST\_CMD\_REG}{0x{}002C}\label{regdesc:CMDREG}
\regfield{SDHOST\_START\_CMD}{1}{31}{0}%
\regfield{(reserved)}{1}{30}{0}%
\regfield{SDHOST\_USE\_HOLE}{1}{29}{1}%
\regfield{(reserved)}{1}{28}{0}%
\regfield{(reserved)}{1}{27}{0}%
\regfield{(reserved)}{1}{26}{0}%
\regfield{(reserved)}{1}{25}{0}%
\regfield{(reserved)}{1}{24}{0}%
\regfield{SDHOST\_CCS\_EXPECTED}{1}{23}{0}%
\regfield{SDHOST\_READ\_CEATA\_DEVICE}{1}{22}{0}%
\regfield{SDHOST\_UPDATE\_CLOCK\_REGISTERS\_ONLY}{1}{21}{0}%
\regfield{SDHOST\_CARD\_NUMBER}{5}{16}{{0x{}00}}%
\regfield{SDHOST\_SEND\_INITIALIZATION}{1}{15}{0}%
\regfield{SDHOST\_STOP\_ABORT\_CMD}{1}{14}{0}%
\regfield{SDHOST\_WAIT\_PRVDATA\_COMPLETE}{1}{13}{0}%
\regfield{SDHOST\_SEND\_AUTO\_STOP}{1}{12}{0}%
\regfield{SDHOST\_TRANSFER\_MODE}{1}{11}{0}%
\regfield{SDHOST\_READ\_WRITE}{1}{10}{0}%
\regfield{SDHOST\_DATA\_EXPECTED}{1}{9}{0}%
\regfield{SDHOST\_CHECK\_RESPONSE\_CRC}{1}{8}{0}%
\regfield{SDHOST\_RESPONSE\_LENGTH}{1}{7}{0}%
\regfield{SDHOST\_RESPONSE\_EXPECT}{1}{6}{0}%
\regfield{SDHOST\_CMD\_INDEX}{6}{0}{{0x{}00}}%
\reglabel{Reset}\regnewline%
\begin{regdesc}\begin{reglist}
\item [SDHOST\_START\_CMD] 开始发送命令。一旦 CIU 响应命令，此位自动清零。当此位置为 1 时，主机不应尝试写任何命令寄存器。如果尝试写寄存器，原始中断寄存器的硬件锁定错误位将会被置为 1。一旦命令被发送并且接收到 SD\_MMC\_CEATA 卡的响应，则原始中断寄存器的 Command Done 位将会被置为 1。(R/W)
\item [SDHOST\_USE\_HOLE]Use Hold 寄存器。(R/W)\\
0：发送给卡的 CMD 和 DATA 旁路 Hold 寄存器；\\
1：发送给卡的 CMD 和 DATA 经过 Hold 寄存器。
\item [SDHOST\_CCS\_EXPECTED] 预期命令完成信号 (CCS) 的配置。(R/W)\\
0：CE-ATA 设备的中断不使能（ATA 控制寄存器中 nIEN = 1）；或者命令不等待设备端的 CCS 信号；\\
1：CE-ATA 设备的中断使能 (nIEN = 0)，并且 RW\_BLK 命令等待 CE-ATA 设备的 CCS 信号。\\
如果命令等待 CE-ATA 设备的 CCS 信号，软件应该将此位置为 1。SD/MMC 置位 RINTSTS 寄存器里的 Data Transfer Over (DTO) 位并且如果 DTO 中断没有被屏蔽，会产生对主机的中断。
\item [SDHOST\_READ\_CEATA\_DEVICE] 读操作标志位。(R/W)\\
0：主机不进行对于 CE-ATA 设备的读操作（RW\_REG 或 RW\_BLK）；\\
1：主机进行对于 CE-ATA 设备的读操作（RW\_REG 或 RW\_BLK）。\\
软件应将此位置为 1 来表明 CE-ATA 设备正在被访问用于读传输。此位用于在执行 CE-ATA 读传输时禁能读数据超时指示。I/O 传输延迟的最大值至少为 10 秒。SD/MMC 在等待来自 CE-ATA 设备的数据时不应指示读取数据超时。
\item [见下页...]
\end{reglist}\end{regdesc}
\end{register}
\addtocounter{Regfloat}{-1}
\begin{register}{H}{SDHOST\_CMD\_REG}{0x{}002C}
\begin{regdesc}\begin{reglist}
\item [接上页...]
\item [SDHOST\_UPDATE\_CLOCK\_REGISTERS\_ONLY] 0：正常命令序列；1：不发送命令，仅更新时钟寄存器的值到卡时钟域内。(R/W)\\
以下寄存器的值被传输到卡时钟域内：CLKDIV、 CLRSRC 和 CLKENA。可改变卡时钟（改变时钟频率、设置是否截断、以及设置低频模式）。
这是为了改变时钟频率或停止时钟，而不必发送命令给卡。\\
在正常命令序列下，当 SDHOST\_UPDATE\_CLOCK\_REGISTERS\_ONLY = 0，以下控制寄存器从 BIU 传输到 CIU：CMD、CMDARG、TMOUT、CTYPE、BLKSIZ 和 BYTCNT。CIU 为新的命令序列使用新的寄存器的值。当此位置为 1 时，由于没有命令被发送给 SD\_MMC\_CEATA 卡，所以没有 Command Done 中断。
\item [SDHOST\_CARD\_NUMBER] 使用中的卡号，表示正在访问的卡的物理插槽编号。 在仅 SD 模式下，支持 2 张卡。(R/W)
\item [SDHOST\_SEND\_INITIALIZATION] 0：在发送命令前不发送初始化序列（80 个时钟周期的 1）；1：在发送命令前发送初始化序列。(R/W)\\
上电后，发送任何命令到卡之前，必须向卡发送 80 个时钟进行初始化。在向卡发送第一个命令时应该将此位置为 1，以便控制器在向卡发送命令之前初始化时钟。
\item [SDHOST\_STOP\_ABORT\_CMD] 0：停止或中止命令，停止命令和中止命令都不会停止当前的数据传输。 如果中止命令发送到当前选择的功能号或不在数据传输模式，则该位应设置为 0；1：停止或中止命令，用于停止当前的数据传输。(R/W)\\ 当开放式预定义数据传输正在进行时，并且主机发出停止或中止命令以停止数据传输时，应将此位置为 1，以使 CIU 的命令/数据状态机可以正确返回到空闲状态。
\item [SDHOST\_WAIT\_PRVDATA\_COMPLETE] 0：即使以前的数据传输尚未完成，也立即发送命令；1：等待前面的数据传输完成后再发送命令。(R/W)\\
SDHOST\_WAIT\_PRVDATA\_COMPLETE = 0 选项通常用来在数据传输时询问卡的状态或停止当前的数据传输。SDHOST\_CARD\_NUMBER 应该与上一个命令相同。
\item [SDHOST\_SEND\_AUTO\_STOP] 0：在数据传输结束时不发送停止命令；1：在数据传送结束时发送停止命令。(R/W)
\item [见下页...]
\end{reglist}\end{regdesc}
\end{register}

\addtocounter{Regfloat}{-1}
\begin{register}{H}{SDHOST\_CMD\_REG}{0x{}002C}
\begin{regdesc}\begin{reglist}
\item [接上页...]
\item [SDHOST\_TRANSFER\_MODE] 0：模块数据传输命令；1: 流数据传输命令。(R/W)\\如果不等待数据则为无关项。
\item [SDHOST\_READ\_WRITE] 0：读卡；1：写卡。(R/W)\\
如果不等待数据则为无关项。
\item [SDHOST\_DATA\_EXPECTED] 0：不等待数据传输；1：等待数据传输。(R/W)
\item [SDHOST\_CHECK\_RESPONSE\_CRC] 0：不检查；1：检查响应 CRC。(R/W)\\
有些命令响应不会返回有效的 CRC 位。软件应禁能对于这些命令的 CRC 检查以便禁能控制器进行 CRC 检查。
\item [SDHOST\_RESPONSE\_LENGTH] 0：等待卡的短响应；1：等待卡的长响应。(R/W)
\item [SDHOST\_RESPONSE\_EXPECT] 0：不等待卡的响应；1：等待卡的响应。(R/W)
\item [SDHOST\_CMD\_INDEX] 命令指数。(R/W)
\end{reglist}\end{regdesc}
\end{register}

\begin{register}{H}{SDHOST\_RESP0\_REG}{0x{}0030}\label{regdesc:RESP0REG}
\regfield{}{32}{0}{{0x{}00000000}}%
\reglabel{Reset}\regnewline%
\begin{regdesc}\begin{reglist}
\item [SDHOST\_RESP0\_REG] 响应的 bit[31:0]。(RO)
\end{reglist}\end{regdesc}
\end{register}


\begin{register}{H}{SDHOST\_RESP1\_REG}{0x{}0034}\label{regdesc:RESP1REG}
\regfield{}{32}{0}{{0x{}00000000}}%
\reglabel{Reset}\regnewline%
\begin{regdesc}\begin{reglist}
\item [SDHOST\_RESP1\_REG] 长响应的 bit[63:32]。(RO)
\end{reglist}\end{regdesc}
\end{register}


\begin{register}{H}{SDHOST\_RESP2\_REG}{0x{}0038}\label{regdesc:RESP2REG}
\regfield{}{32}{0}{{0x{}00000000}}%
\reglabel{Reset}\regnewline%
\begin{regdesc}\begin{reglist}
\item [SDHOST\_RESP2\_REG] 长响应的 bit[95:64]。(RO)
\end{reglist}\end{regdesc}
\end{register}


\begin{register}{H}{SDHOST\_RESP3\_REG}{0x{}003C}\label{regdesc:RESP3REG}
\regfield{}{32}{0}{{0x{}00000000}}%
\reglabel{Reset}\regnewline%
\begin{regdesc}\begin{reglist}
\item [SDHOST\_RESP3\_REG] 长响应的 bit[127:96]。(RO)
\end{reglist}\end{regdesc}
\end{register}


\begin{register}{H}{SDHOST\_MINTSTS\_REG}{0x{}0040}\label{regdesc:MINTSTSREG}
\regfield{(reserved)}{14}{18}{{0x{}0000}}%
\regfield{SDHOST\_SDIO\_INTERRUPT\_MSK}{2}{16}{{0x{}0}}%
\regfield{SDHOST\_INT\_STATUS\_MSK}{16}{0}{{0x{}0000}}%
\reglabel{Reset}\regnewline%
\begin{regdesc}\begin{reglist}
\item [SDHOST\_SDIO\_INTERRUPT\_MSK]SDIO 中断的屏蔽位，每个卡占一个位。Bit[17:16] 分别对应卡 1 和卡 0。只有对应的 sdhost\_sdio\_int\_mask 位被置为 1 时，SDIO 中断才会使能（置位屏蔽位使能中断）。(RO)
\item [SDHOST\_INT\_STATUS\_MSK] 只有当中断屏蔽寄存器中的对应位被置为 1 时，中断才会使能。(RO)\\
Bit 15 (EBE): 结束位错误／无 CRC 错误；\\
Bit 14 (ACD): 自动命令结束；\\
Bit 13 (SBE/BCI): 接收启动位错误；\\
Bit 12 (HLE): 硬件锁定写入错误\\
Bit 11 (FRUN): FIFO 空/满 错误；\\
Bit 10 (HTO): 主机填充数据超时；\\
Bit 9 (DRTO): 数据读取超时；\\
Bit 8 (RTO): 响应超时； \\
Bit 7 (DCRC): 数据 CRC 错误； \\
Bit 6 (RCRC): 响应 CRC 错误； \\
Bit 5 (RXDR): 接收 FIFO 数据请求； \\
Bit 4 (TXDR): 发送 FIFO 数据请求； \\
Bit 3 (DTO): 数据传输结束； \\
Bit 2 (CD): 命令执行完毕； \\
Bit 1 (RE): 响应错误；\\
Bit 0 (CD): 卡检测。
\end{reglist}\end{regdesc}
\end{register}


\begin{register}{H}{SDHOST\_RINTSTS\_REG}{0x{}0044}\label{regdesc:RINTSTSREG}
\regfield{(reserved)}{14}{18}{{0x{}0000}}%
\regfield{SDHOST\_SDIO\_INTERRUPT\_RAW}{2}{16}{{0x{}0}}%
\regfield{SDHOST\_INT\_STATUS\_RAW}{16}{0}{{0x{}0000}}%
\reglabel{Reset}\regnewline%
\begin{regdesc}\begin{reglist}
\item [SDHOST\_SDIO\_INTERRUPT\_RAW] 来自 SDIO 卡的中断，一个卡占一个位。Bit[17:16] 分别对应卡 1 和卡 0。置位某位就把相应的中断位清零，写 0 无效。(R/W)\\
0：没有来自卡的 SDIO 中断；\\
1：有来自卡的 SDIO 中断。
\item [SDHOST\_INT\_STATUS\_RAW] 置位某位就把相应的中断位清零，写 0 无效。无论中断屏蔽状态如何，这些中断位都会被记录。(R/W)\\
Bit 15 (EBE): 结束位错误／无 CRC 错误；\\
Bit 14 (ACD): 自动命令结束；\\
Bit 13 (SBE/BCI): 接收启动为错误；\\
Bit 12 (HLE): 硬件锁定写入错误\\
Bit 11 (FRUN): FIFO 空/满 错误；\\
Bit 10 (HTO): 主机填充数据超时；\\
Bit 9 (DRTO): 数据读取超时；\\
Bit 8 (RTO): 响应超时； \\
Bit 7 (DCRC): 数据 CRC 错误； \\
Bit 6 (RCRC): 响应 CRC 错误； \\
Bit 5 (RXDR): 接收 FIFO 数据请求； \\
Bit 4 (TXDR): 发送 FIFO 数据请求； \\
Bit 3 (DTO): 数据传输结束； \\
Bit 2 (CD): 命令执行完毕； \\
Bit 1 (RE): 响应错误；\\
Bit 0 (CD): 卡检测。
\end{reglist}\end{regdesc}
\end{register}


\begin{register}{H}{SDHOST\_STATUS\_REG}{0x{}0048}\label{regdesc:STATUSREG}
\regfield{(reserved)}{1}{31}{0}%
\regfield{(reserved)}{1}{30}{0}%
\regfield{SDHOST\_FIFO\_COUNT}{13}{17}{{0x{}000}}%
\regfield{SDHOST\_RESPONSE\_INDEX}{6}{11}{{0x{}00}}%
\regfield{SDHOST\_DATA\_STATE\_MC\_BUSY}{1}{10}{1}%
\regfield{SDHOST\_DATA\_BUSY}{1}{9}{1}%
\regfield{SDHOST\_DATA\_3\_STATUS}{1}{8}{1}%
\regfield{SDHOST\_COMMAND\_FSM\_STATES}{4}{4}{{0x{}1}}%
\regfield{SDHOST\_FIFO\_FULL}{1}{3}{0}%
\regfield{SDHOST\_FIFO\_EMPTY}{1}{2}{1}%
\regfield{SDHOST\_FIFO\_TX\_WATERMARK}{1}{1}{1}%
\regfield{SDHOST\_FIFO\_RX\_WATERMARK}{1}{0}{0}%
\reglabel{Reset}\regnewline%
\begin{regdesc}\begin{reglist}
\item [SDHOST\_FIFO\_COUNT]FIFO 计数器，FIFO 中被填充的地址的数量。(RO)
\item [SDHOST\_RESPONSE\_INDEX] 前一个响应的指数，包括内核发送的任何自动停止的响应。(RO)
\item [SDHOST\_DATA\_STATE\_MC\_BUSY] 数据发送或接收状态机忙。(RO)
\item [SDHOST\_DATA\_BUSY] 数据线 sdhost\_card\_data[0] 的值取反。(RO)\\
0：卡数据不忙；\\
1：卡数据忙。
\item [SDHOST\_DATA\_3\_STATUS] 数据线 sdhost\_card\_data[3] 上的值，检查卡是否存在。(RO)\\
0：卡不存在；\\
1：卡存在。
\item [SDHOST\_COMMAND\_FSM\_STATES] 命令状态机状态。(RO)\\
0: 空闲；\\
1: 发送初始序列； \\
2: 发送命令开始位； \\
3: 发送命令发送位；\\
4: 发送命令指数和参数；\\
5: 发送命令 CRC7；\\
6: 发送命令结束位；\\
7: 接收响应开始位；\\
8: 接收响应 IRQ 响应；\\
9: 接收响应发送位；\\
10: 接收响应命令指数；\\
11: 接收响应数据；\\
12: 接收响应 CRC7；\\
13: 接收响应结束位；\\
14: 命令路径等待 NCC；\\
15: 等待，命令-响应回转。
\item [SDHOST\_FIFO\_FULL]FIFO 为满的状态。(RO)
\item [SDHOST\_FIFO\_EMPTY]FIFO 为空的状态。(RO)
\item [SDHOST\_FIFO\_TX\_WATERMARK]FIFO 达到发送阈值，不是数据传输的必要条件。(RO)
\item [SDHOST\_FIFO\_RX\_WATERMARK]FIFO 达到接收阈值，不是数据传输的必要条件。(RO)
\end{reglist}\end{regdesc}
\end{register}


\begin{register}{H}{SDHOST\_FIFOTH\_REG}{0x{}004C}\label{regdesc:FIFOTHREG}
\regfield{(reserved)}{1}{31}{0}%
\regfield{SDHOST\_DMA\_MULTIPLE\_TRANSACTION\_SIZE}{3}{28}{{0x{}0}}%
\regfield{(reserved)}{1}{27}{0}%
\regfield{SDHOST\_RX\_WMARK}{11}{16}{xxxxxxxxxxx}%
\regfield{(reserved)}{4}{12}{0000}%
\regfield{SDHOST\_TX\_WMARK}{12}{0}{{0x{}000}}%
\reglabel{Reset}\regnewline%
\begin{regdesc}\begin{reglist}
\item [SDHOST\_DMA\_MULTIPLE\_TRANSACTION\_SIZE] 突发传输的字节数，配置的值应与 DMA 控制器多个传输大小 SDHOST\_SRC/DEST\_MSIZE 相同。000: 1 字节传输；001: 4 字节传输；010: 8 字节传输；011: 16 字节传输；100: 32 字节传输；101: 64 字节传输；110: 128 字节传输；111: 256 字节传输。(R/W)
\item [SDHOST\_RX\_WMARK] 当接收数据给卡时 FIFO 的阈值。当 FIFO 数据计数大于该数值时，DMA/FIFO 请求被提出。 在数据包结束期间，无论阈值大小如何，都会生成请求，以完成剩余的数据传输。在非 DMA 模式下，当接收 FIFO 阈值 (RXDR) 中断使能时，则会产生中断，而不是 DMA 请求。 如果阈值大于任何剩余数据，则不产生中断。当主机看见数据发送结束中断时，主机应该读取剩下的数据。在 DMA 模式下，在数据包结束时，即使剩余的字节数少于阈值，DMA  请求也会在数据传输结束中断设置之前进行单次传输以清除所有剩余的字节。(R/W)
\item [SDHOST\_TX\_WMARK] 当发送数据给卡时 FIFO 的阈值。当 FIFO 数据计数小于等于该数值时，DMA/FIFO 请求被提出。如果使能中断，则中断发生。 在数据包结束期间，无论阈值大小如何，都会生成请求。在非 DMA 模式下，当发送 FIFO 阈值 (TXDR) 中断使能时，则会产生中断，而不是 DMA 请求。在数据包结束期间，在最后一个中断时，主机负责仅用所需的剩余字节填充 FIFO（不是在 FIFO 满之前或在 CIU 完成数据传输之后，因为 FIFO 可能不为空）。 在 DMA 模式下，在数据包结束时，如果最后一次传输比突发传输小，DMA 控制器将执行单个周期，直到传输所需的字节。(R/W)
\end{reglist}\end{regdesc}
\end{register}


\begin{register}{H}{SDHOST\_CDETECT\_REG}{0x{}0050}\label{regdesc:CDETECTREG}
\regfield{(reserved)}{30}{2}{{0x{}0000000}}%
\regfield{SDHOST\_CARD\_DETECT\_N}{2}{0}{{0x{}0}}%
\reglabel{Reset}\regnewline%
\begin{regdesc}\begin{reglist}
\item [SDHOST\_CARD\_DETECT\_N] 信号线 sdhost\_card\_detect\_n 输入端口（每个卡一个位）的值。0 代表卡存在。只有 相应位被执行。(RO)
\end{reglist}\end{regdesc}
\end{register}


\begin{register}{H}{SDHOST\_WRTPRT\_REG}{0x{}0054}\label{regdesc:WRTPRTREG}
\regfield{(reserved)}{30}{2}{{0x{}0000000}}%
\regfield{SDHOST\_WRITE\_PROTECT}{2}{0}{{0x{}0}}%
\reglabel{Reset}\regnewline%
\begin{regdesc}\begin{reglist}
\item [SDHOST\_WRITE\_PROTECT] 信号线 sdhost\_card\_write\_prt 输入端口（每个卡一个位）的值。1 表示写保护。只有 的相应位被执行。(RO)
\end{reglist}\end{regdesc}
\end{register}


\begin{register}{H}{SDHOST\_TCBCNT\_REG}{0x{}005C}\label{regdesc:TCBCNTREG}
\regfield{}{32}{0}{{0x{}00000000}}%
\reglabel{Reset}\regnewline%
\begin{regdesc}\begin{reglist}
\item [SDHOST\_TCBCNT\_REG]CIU 模块以及传输给卡的字节数。(RO)
\end{reglist}\end{regdesc}
\end{register}


\begin{register}{H}{SDHOST\_TBBCNT\_REG}{0x{}0060}\label{regdesc:TBBCNTREG}
\regfield{}{32}{0}{{0x{}00000000}}%
\reglabel{Reset}\regnewline%
\begin{regdesc}\begin{reglist}
\item [SDHOST\_TBBCNT\_REG] 主机/DMA 和 BIU FIFO 之间传输的字节数。(RO)
\end{reglist}\end{regdesc}
\end{register}


\begin{register}{H}{SDHOST\_DEBNCE\_REG}{0x{}0064}\label{regdesc:DEBNCEREG}
\regfield{(reserved)}{8}{24}{00000000}%
\regfield{SDHOST\_DEBOUNCE\_COUNT}{24}{0}{{0x{}000000}}%
\reglabel{Reset}\regnewline%
\begin{regdesc}\begin{reglist}
\item [SDHOST\_DEBOUNCE\_COUNT] 抖动消除滤波器逻辑使用的主机时钟数。典型的去抖动时间为 5 \verb+~+ 25 ms，防止插卡或移除卡的时候的不稳定性。(R/W)
\end{reglist}\end{regdesc}
\end{register}


\begin{register}{H}{SDHOST\_USRID\_REG}{0x{}0068}\label{regdesc:USRIDREG}
\regfield{}{32}{0}{{0x{}00000000}}%
\reglabel{Reset}\regnewline%
\begin{regdesc}\begin{reglist}
\item [SDHOST\_USRID\_REG] 用户识别寄存器。此寄存器也可以作为用户的寄存器使用。(R/W)
\end{reglist}\end{regdesc}
\end{register}


\begin{register}{H}{SDHOST\_VERID\_REG}{0x{}006C}\label{regdesc:VERIDREG}
\regfield{}{32}{0}{{0x{}5432270A}}%
\reglabel{Reset}\regnewline%
\begin{regdesc}\begin{reglist}
\item [SDHOST\_VERSIONID\_REG] 表示硬件版本。固件也可以读取该寄存器。 (RO)
\end{reglist}\end{regdesc}
\end{register}


\begin{register}{H}{SDHOST\_HCON\_REG}{0x{}0070}\label{regdesc:HCONREG}
\regfield{(reserved)}{5}{27}{{0x{}0}}%
\regfield{(reserved)}{1}{26}{{0x{}0}}%
\regfield{SDHOST\_NUM\_CLK\_DIV\_REG}{2}{24}{{0x{}3}}%
\regfield{(resvrved)}{1}{23}{{0x{}1}}%
\regfield{SDHOST\_HOLD\_REG}{1}{22}{{0x{}1}}%
\regfield{SDHOST\_RAM\_INDISE\_REG}{1}{21}{{0x{}0}}%
\regfield{SDHOST\_DMA\_WIDTH\_REG}{3}{18}{{0x{}1}}%
\regfield{(reserved)}{2}{16}{{0x{}0}}%
\regfield{SDHOST\_ADDR\_WIDTH\_REG}{6}{10}{{0x{}13}}%
\regfield{SDHOST\_DATA\_WIDTH\_REG}{3}{7}{{0x{}1}}%
\regfield{SDHOST\_BUS\_TYPE\_REG}{1}{6}{{0x{}1}}%
\regfield{SDHOST\_CARD\_NUM\_REG}{5}{1}{{0x{}1}}%
\regfield{SDHOST\_CARD\_TYPE\_REG}{1}{0}{{0x{}1}}%
\reglabel{Reset}\regnewline%
\begin{regdesc}\begin{reglist}
\item [SDHOST\_NUM\_CLK\_DIV\_REG] 表示设计中有 4 个时钟分频器。(RO)
\item [SDHOST\_HOLD\_REG] 表示数据路径中有一个保持寄存器。(RO)
\item [SDHOST\_RAM\_INDISE\_REG] 表示 SDMMC 模块内部包含 RAM。(RO)
\item [SDHOST\_DMA\_WIDTH\_REG] 表示 DMA 数据宽度为 32。(RO)
\item [SDHOST\_ADDR\_WIDTH\_REG] 表示寄存器地址宽度为 32。(RO)
\item [SDHOST\_DATA\_WIDTH\_REG] 表示寄存器数据宽度为 32。(RO)
\item [SDHOST\_BUS\_TYPE\_REG] 表示寄存器配置为 APB 总线。(RO)
\item [SDHOST\_CARD\_NUM\_REG] 表示支持的卡数量为 2。(RO)
\item [SDHOST\_CARD\_TYPE\_REG] 表示硬件支持 SDIO 和 MMC。(RO)
\end{reglist}\end{regdesc}
\end{register}


\begin{register}{H}{SDHOST\_UHS\_REG}{0x{}0074}\label{regdesc:UHSREG}\regfield{reserved}{14}{18}{{0x{}0000}}%
\regfield{SDHOST\_DDR\_REG}{2}{16}{{0x{}0}}%
\regfield{reserved}{16}{0}{{0x{}0000}}%
\reglabel{Reset}\regnewline%
\begin{regdesc}\begin{reglist}
\item [SDHOST\_DDR\_REG]配置卡是否使用 DDR（双倍数据速率）模式。每张卡一位，位 0 对应卡 0。\\
0：非 DDR 模式\\
1：DDR 模式\\ (R/W)
\end{reglist}\end{regdesc}
\end{register}


\begin{register}{H}{SDHOST\_RST\_N\_REG}{0x{}0078}\label{regdesc:RSTNREG}
\regfield{(reserved)}{30}{2}{{0x{}00000000}}%
\regfield{SDHOST\_RST\_CARD\_RESET}{2}{0}{{0x{}1}}%
\reglabel{Reset}\regnewline%
\begin{regdesc}\begin{reglist}
\item [SDHOST\_RST\_CARD\_RESET] 硬件复位。(R/W)\\
1：工作模式；\\
0：复位。\\
这些位让卡进入空闲状态，主机工作时需要被重新初始化。SDHOST\_RST\_CARD\_RESET[0] 应被置为 1'b0 来复位卡 0。SDHOST\_RST\_CARD\_RESET[1] 应被置为 1'b0 来复位卡 1。
\end{reglist}\end{regdesc}
\end{register}


\begin{register}{H}{SDHOST\_BMOD\_REG}{0x{}0080}\label{regdesc:BMODREG}
\regfield{(reserved)}{21}{11}{000000000000000000000}%
\regfield{SDHOST\_BMOD\_PBL}{3}{8}{{0x{}0}}%
\regfield{SDHOST\_BMOD\_DE}{1}{7}{0}%
\regfield{(reserved)}{5}{2}{{0x{}00}}%
\regfield{SDHOST\_BMOD\_FB}{1}{1}{0}%
\regfield{SDHOST\_BMOD\_SWR}{1}{0}{0}%
\reglabel{Reset}\regnewline%
\begin{regdesc}\begin{reglist}
\item [SDHOST\_BMOD\_PBL] 可编程突发长度。 这些位指示在一个 IDMAC (Internal DMA Control) 传输中要执行的最大节拍数。 IDMAC 每次在主机总线上开始突发传输时将总是尝试在 PBL 中指定的突发值。允许的值为 1、4、8、16、32、64、128 和 256。该值是 FIFOTH 寄存器的 MSIZE 的镜像。 要修改此值，将所需的值写入 FIFOTH 寄存器。 这是一个编码值，如下所示：(RO)\\
000: 1 字节传输；\\
001: 4 字节传输；\\
010: 8 字节传输；\\
011: 16 字节传输；\\
100: 32 字节传输；\\
101: 64 字节传输；\\
110: 128 字节传输；\\
111: 256 字节传输。\\
PBL 是只读值，并且仅适用于数据访问，不适用于链表访问。
\item [SDHOST\_BMOD\_DE]IDMAC 使能位。置位后 IDMAC 使能。(RO)
\item [SDHOST\_BMOD\_FB] 固定突发。 控制 AHB 主接口是否执行固定突发传输。 置位时，AHB 将在正常突发传输开始期间仅使用 SINGLE、INCR4、INCR8 或 INCR16。 当复位时，AHB 将使用 SINGLE 和 INCR 突发传输操作。(R/W)
\item [SDHOST\_BMOD\_SWR] 软件复位。当置位时，DMA 控制器复位所有内部寄存器。一个时钟周期后自动清零。(R/W)
\end{reglist}\end{regdesc}
\end{register}


\begin{register}{H}{SDHOST\_PLDMND\_REG}{0x{}0080}\label{regdesc:PLDMNDREG}
\regfield{}{32}{0}{{0x{}00000000}}%
\reglabel{Reset}\regnewline%
\begin{regdesc}\begin{reglist}
\item [SDHOST\_PLDMND\_REG] 轮询需求。 如果链表的 OWNER 位未置位，则 状态机 进入挂起状态。主机需要对这个寄存器写入任意值，以使 IDMAC 状态机 恢复正常读取链表的操作。 这是一个只写寄存器。(WO)
\end{reglist}\end{regdesc}
\end{register}


\begin{register}{H}{SDHOST\_DBADDR\_REG}{0x{}0088}\label{regdesc:DBADDRREG}
\regfield{}{32}{0}{{0x{}00000000}}%
\reglabel{Reset}\regnewline%
\begin{regdesc}\begin{reglist}
\item [SDHOST\_DBADDR\_REG] 链表的开始。 包含第一个链表的基址。最低效位 (LSB bit) [1:0] 被忽略，并由 IDMAC 内部全部取为零，因此这些 LSB 位可被视为只读。(R/W)
\end{reglist}\end{regdesc}
\end{register}


\begin{register}{H}{SDHOST\_IDSTS\_REG}{0x{}008C}\label{regdesc:IDSTSREG}
\regfield{(reserved)}{15}{17}{000000000000000}%
\regfield{SDHOST\_IDSTS\_FSM}{4}{13}{{0x{}0}}%
\regfield{SDHOST\_IDSTS\_FBE\_CODE}{3}{10}{{0x{}0}}%
\regfield{SDHOST\_IDSTS\_AIS}{1}{9}{0}%
\regfield{SDHOST\_IDSTS\_NIS}{1}{8}{0}%
\regfield{(reserved)}{2}{6}{00}%
\regfield{SDHOST\_IDSTS\_CES}{1}{5}{0}%
\regfield{SDHOST\_IDSTS\_DU}{1}{4}{0}%
\regfield{(reserved)}{1}{3}{0}%
\regfield{SDHOST\_IDSTS\_FBE}{1}{2}{0}%
\regfield{SDHOST\_IDSTS\_RI}{1}{1}{0}%
\regfield{SDHOST\_IDSTS\_TI}{1}{0}{0}%
\reglabel{Reset}\regnewline%
\begin{regdesc}\begin{reglist}
\item [SDHOST\_IDSTS\_FSM]DMAC 状态机 当前状态。(RO)\\
0: DMA\_IDLE（状态机空闲状态）；\\
1: DMA\_SUSPEND（状态机暂停状态）；\\
2: DESC\_RD（读链表状态）；\\
3: DESC\_CHK（检查链表状态）；\\
4: DMA\_RD\_REQ\_WAIT（状态机读数据请求等待状态）；\\
5: DMA\_WR\_REQ\_WAIT（状态机写数据请求等待状态）；\\
6: DMA\_RD（状态机读数据状态）； \\
7: DMA\_WR（状态机写数据状态）； \\
8: DESC\_CLOSE（当前链表使用完关闭状态）。
\item [SDHOST\_IDSTS\_FBE\_CODE] 致命总线错误代码。 表明导致总线错误的错误类型。仅当设置致命总线错误位 IDSTS[2] 被置位时有效。 此字段不生成中断。(RO)\\
001：发送期间接收到主机中止发送 ；\\
010：接收期间接收到主机中止接收；\\
其他：保留。
\item [SDHOST\_IDSTS\_AIS] 异常中断汇总。以下各项的逻辑或：(R/W)\\
IDSTS[2]：致命总线中断；\\
IDSTS[4]：SDHOST\_IDSTS\_DU 位中断。\\
只有未屏蔽的位影响该位。 这是一个粘滞位，必须在每次清零可以引起 AIS 置 1 的相应位时清零。写 1 清零该位。
\item [SDHOST\_IDSTS\_NIS] 正常中断汇总。以下各项的逻辑或：(R/W)\\
IDSTS[0]：发送中断；\\
IDSTS[1]：接收中断。\\
只有未屏蔽的位影响该位。 这是一个粘滞位，必须在每次清零可以引起 NIS 置 1 的相应位时清零。写 1 清零该位。
\item [见下页...]
\end{reglist}\end{regdesc}
\end{register}

\addtocounter{Regfloat}{-1}
\begin{register}{H}{SDHOST\_IDSTS\_REG}{0x{}008C}
\begin{regdesc}\begin{reglist}
\item [接上页...]
\item [SDHOST\_IDSTS\_CES] 卡错误汇总。指示发送/接收卡的传输状态，也出现在 RINTSTS 中。 表示以下位的逻辑或：(R/W)\\
EBE：结束位错误；\\
RTO：响应超时/引导确认超时错误；\\
RCRC：响应 CRC 错误；\\
SBE：启动位错误；\\
DRTO：数据读取超时/BDS 超时错误；\\
DCRC：接收的数据 CRC 错误；\\
RE：响应错误。\\
写 1 清零该位。IDMAC 的中止条件取决于此 CES 位的配置。 如果 CES 位被使能，则 IDMAC 在响应错误时中止。
\item [SDHOST\_IDSTS\_DU] 链表不可用中断。当链表由于 OWNER 位= 0 (DES0 [31] = 0) 而不可用时，该位置 1。 写 1 清零该位。(R/W)
\item [SDHOST\_IDSTS\_FBE] 致命总线错误中断。表示发生总线错误 (IDSTS[12:10])。当该位置 1 时，DMA 禁止所有总线访问。写 1 清零该位。(R/W)
\item [SDHOST\_IDSTS\_RI] 接收中断。表示链表的数据接收完成。写 1 清零该位。(R/W)
\item [SDHOST\_IDSTS\_TI] 发送中断。表示链表的数据发送完成。写 1 清零该位。(R/W)
\end{reglist}\end{regdesc}
\end{register}


\begin{register}{H}{SDHOST\_IDINTEN\_REG}{0x{}0090}\label{regdesc:IDINTENREG}
\regfield{(reserved)}{22}{10}{0000000000000000000000}%
\regfield{SDHOST\_IDINTEN\_AI}{1}{9}{0}%
\regfield{SDHOST\_IDINTEN\_NI}{1}{8}{0}%
\regfield{(reserved)}{2}{6}{00}%
\regfield{SDHOST\_IDINTEN\_CES}{1}{5}{0}%
\regfield{SDHOST\_IDINTEN\_DU}{1}{4}{0}%
\regfield{(reserved)}{1}{3}{0}%
\regfield{SDHOST\_IDINTEN\_FBE}{1}{2}{0}%
\regfield{SDHOST\_IDINTEN\_RI}{1}{1}{0}%
\regfield{SDHOST\_IDINTEN\_TI}{1}{0}{0}%
\reglabel{Reset}\regnewline%
\begin{regdesc}\begin{reglist}
\item [SDHOST\_IDINTEN\_AI] 异常中断汇总使能位。置 1 时，使能异常中断。该位使能以下位：(R/W)\\
IDINTEN [2]：致命总线错误中断；\\
IDINTEN [4]：DU 中断。
\item [SDHOST\_IDINTEN\_NI] 正常中断汇总使能位。置 1 时，使能正常中断。重置时，禁能正常中断。该位使能以下位：(R/W)\\
IDINTEN[0]：发送中断；\\
IDINTEN[1]：接收中断。
\item [SDHOST\_IDINTEN\_CES] 卡错误汇总中断使能位。置 1 时使能卡中断汇总。(R/W)
\item [SDHOST\_IDINTEN\_DU] 链表不可用中断。当与异常中断汇总使能位一起设置时，将使能 DU 中断。(R/W)
\item [SDHOST\_IDINTEN\_FBE] 致命总线错误使能位。当与异常中断汇总使能位一起设置时，将使能致命总线错误中断。复位时，致命总线错误使能中断被禁能。(R/W)
\item [SDHOST\_IDINTEN\_RI] 接收中断使能位。当与正常中断汇总使能位一起设置时，将使能接收中断。复位时，接收中断被禁能。(R/W)
\item [SDHOST\_IDINTEN\_TI] 发送中断使能位。当与正常中断汇总使能位一起设置时，将使能发送中断。复位时，发送中断被禁能。(R/W)
\end{reglist}\end{regdesc}
\end{register}


\begin{register}{H}{SDHOST\_DSCADDR\_REG}{0x{}0094}\label{regdesc:DSCADDRREG}
\regfield{}{32}{0}{{0x{}00000000}}%
\reglabel{Reset}\regnewline%
\begin{regdesc}\begin{reglist}
\item [SDHOST\_DSCADDR\_REG] 主机链表地址指针，在操作期间由 IDMAC 更新，并在复位时清零。该寄存器指向由 IDMAC 读取的当前链表的起始地址。(RO)
\end{reglist}\end{regdesc}
\end{register}


\begin{register}{H}{SDHOST\_BUFADDR\_REG}{0x{}0098}\label{regdesc:BUFADDRREG}
\regfield{}{32}{0}{{0x{}00000000}}%
\reglabel{Reset}\regnewline%
\begin{regdesc}\begin{reglist}
\item [SDHOST\_BUFADDR\_REG] 主机缓冲区地址指针，在操作期间由 IDMAC 更新，并在复位时清零。该寄存器指向由 IDMAC 访问的当前数据缓冲区地址。(RO)
\end{reglist}\end{regdesc}
\end{register}


\begin{register}{H}{SDHOST\_CARDTHRCTL\_REG}{0x{}0100}\label{regdesc:CARDTHRCTLREG}
\regfield{SDHOST\_CARDTHRESHOLD\_REG}{16}{16}{{0x{}000}}%
\regfield{(reserved)}{13}{3}{{0x{}00}}%
\regfield{SDHOST\_CARDWRTHREN\_REG}{1}{2}{{0}}%
\regfield{SDHOST\_CARDCLRINTEN\_REG}{1}{1}{{0}}%
\regfield{SDHOST\_CARDRDTHREN\_REG}{1}{0}{{0}}%
\reglabel{Reset}\regnewline%
\begin{regdesc}\begin{reglist}
\item [SDHOST\_CARDTHRESHOLD\_REG] 配置卡读取阈值。内部 FIFO 大小为 512。当 SDHOST\_CARDERTHREN\_REG 被设置为 1 或 SDHOST\_CARDRDTHREN\_REG 被设置为 1 时，该寄存器生效。 (R/W)\\
\item [SDHOST\_CARDWRTHREN\_REG] 在 HS400 模式启用时适用。\\
0：关闭卡写阈值功能。\\
1：启用卡写阈值功能。\\ (R/W)
\item [SDHOST\_CARDCLRINTEN\_REG] 配置是否使能忙碌清除中断生成。\\
0：不使能\\
1：使能\\ (R/W)
\item [SDHOST\_CARDRDTHREN\_REG] 配置是否使能卡读取阈值。\\
0：不使能\\
1：使能\\ (R/W)
\end{reglist}\end{regdesc}
\end{register}


\begin{register}{H}{SDHOST\_EMMC\_DDR\_REG}{0x{}010C}\label{regdesc:EMMCDDRREG}
\regfield{SDHOST\_HS400\_MODE\_REG}{1}{31}{{0x{}0}}%
\regfield{(reserved)}{29}{2}{{0x{}000000}}%
\regfield{SDHOST\_HALFSTARTBIT\_REG}{2}{0}{{0x{}0}}%
\reglabel{Reset}\regnewline%
\begin{regdesc}\begin{reglist}
\item [SDHOST\_HS400\_MODE\_REG] 置 1 启用 HS400 模式。\\
\item [SDHOST\_HALFSTARTBIT\_REG] 配置开始位检测机制的开始位检测长度。每张卡一位，位 0 对应卡 0。eMMC4.5 及以上版本此位需为 1，SD 应用此位需为 0。对于 eMMC4.5，起始位检测模式可配置为：\\
0：完整周期\\
1：少于一个完整周期\\(R/W)
\end{reglist}\end{regdesc}
\end{register}

\begin{register}{H}{SDHOST\_ENSHIFT\_REG}{0x{}0110}\label{regdesc:ENSHIFTREG}
\regfield{(reserved)}{28}{4}{{0x{}0000000}}%
\regfield{SDHOST\_ENABLE\_SHIFT\_REG}{4}{0}{{0x{}0}}%
\reglabel{Reset}\regnewline%
\begin{regdesc}\begin{reglist}
\item [SDHOST\_ENABLE\_SHIFT\_REG] 配置设计中默认的相位偏移量。每张卡 2 位，位 1:0 对应卡 0。\\
0x0：默认相位偏移\\
0x1：偏移到下个上升沿\\
0x2：偏移到下个下降沿\\
0x3：保留\\ (R/W)
\end{reglist}\end{regdesc}
\end{register}


\begin{register}{H}{SDHOST\_BUFFIFO\_REG}{0x{}0200}\label{regdesc:BUFFIFOREG}
\regfield{}{32}{0}{{0x{}00000000}}%
\reglabel{Reset}\regnewline%
\begin{regdesc}\begin{reglist}
\item [SDHOST\_BUFFIFO\_REG]CPU 通过 FIFO 读写发送的数据。该寄存器指向当前数据 FIFO。(RO)
\end{reglist}\end{regdesc}
\end{register}

\begin{register}{H}{SDHOST\_CLK\_DIV\_EDGE\_REG}{0x{}0800}\label{regdesc:CLKEDGESEL}
\regfield{(reserved)}{11}{21}{{0x{}000}}%
\regfield{SDHOST\_CCLKIN\_EDGE\_N}{4}{17}{{0x{}1}}%
\regfield{SDHOST\_CCLKIN\_EDGE\_L}{4}{13}{{0x{}0}}%
\regfield{SDHOST\_CCLKIN\_EDGE\_H}{4}{9}{{0x{}1}}%
\regfield{SDHOST\_CCLKIN\_EDGE\_SLF\_SEL}{3}{6}{{0x{}0}}%
\regfield{SDHOST\_CCLKIN\_EDGE\_SAM\_SEL}{3}{3}{{0x{}0}}%
\regfield{SDHOST\_CCLKIN\_EDGE\_DRV\_SEL}{3}{0}{{0x{}0}}%
\reglabel{Reset}\regnewline%
\begin{regdesc}\begin{reglist}
\item[CCLKIN\_EDGE\_N] 值应与 CCLKIN\_EDGE\_L 相同。(R/W)
\item[CCLKIN\_EDGE\_L] 分频时钟的低电平，值应比 CCLKIN\_EDGE\_H 大。(R/W)
\item[CCLKIN\_EDGE\_H] 分频时钟的高电平，值应比 CCLKIN\_EDGE\_L 小。(R/W)
\item[CCLKIN\_EDGE\_SLF\_SEL] 用于选择内部时钟信号的相位，90 度相位、180 度相位或 270 度相位。(R/W)
\item[CCLKIN\_EDGE\_SAM\_SEL] 用于选择输入时钟信号的相位，90 度相位、180 度相位或 270 度相位。(R/W)
\item[CCLKIN\_EDGE\_DRV\_SEL] 用于选择输出时钟信号的相位，90 度相位、180 度相位或 270 度相位。(R/W)\\
\end{reglist}\end{regdesc}
说明：SD/MMC 使用该寄存器分频 160M 时钟 (CCLKIN\_EDGE\_H/CCLKIN\_EDGE\_L)。输出时钟通过该寄存器和寄存器 SDHOST\_CLKDIV\_REG 连接 SDIO 从机分频器，使用 SDHOST\_CLKSRC\_REG 时可选择 4 个时钟源。


\end{register}
